% 取消下一行注释可生成 handout（无逐步显示）；保持注释则生成正常 slide。
% \PassOptionsToClass{handout}{beamer}
\documentclass[Main-Prob-All.tex]{subfiles}
\begin{document}
\title[概率论]{第二十讲：特征函数}
%\author[张鑫 {\rm Email: xzhangseu@seu.edu.cn} ]{\large 张 鑫}
\institute{\large \textrm{Email: xzhangseu@seu.edu.cn} \\ \quad  \\
	\vspace{0.3cm}
	% \trc{公共邮箱: \textrm{zy.prob@qq.com}\\
		%\hspace{-1.7cm}  密 \qquad 码: \textrm{seu!prob}}
}
\date{}

{ \setbeamertemplate{footline}{}
	\begin{frame}
		\titlepage
	\end{frame}
}



\subsection{特征函数}
\begin{frame}
	\frametitle{复随机变量及其期望}
	\begin{defi}
		若 $X (\omega),Y (\omega)$ 为定义在 $\Omega$ 上的实值随机变量，则称 $Z (\omega)=X (\omega)+iY (\omega)$ 为复随机变量；称 $\overline{Z}=X (\omega)-iY (\omega)$ 为 $Z (\omega)$ 的复共轭随机变量；称 $|Z|:=\sqrt{X^2+Y^2}$ 为复随机变量 $Z$ 的模.
	\end{defi}
	\pause
	\begin{defi}
		若随机变量 $X,Y$ 的期望 $E (X),E (Y)$ 都存在，则复随机变量 $Z$ 的数学期望定义为 $E (Z):=E (X)+iE (Y)$.
	\end{defi}

	\pause
	\begin{itemize}[<+-|alert@+>]
		\item $Z_1=X_1+iY_1$ 与 $Z_2=X_2+iY_2$ 独立当且仅当 $(X_1,Y_1)$ 与 $(X_2,Y_2)$ 独立；
		\item $E(e^{iX})=E(\cos X)+i E(\sin X)$;
		\item $|e^{iX}|=\sqrt{\cos^2 X+\sin^2 X}=1$;
		\item 若 $X,Y$ 独立，则 $e^{iX}$ 与 $e^{iY}$ 也独立.
	\end{itemize}

\end{frame}


\begin{frame}
	\frametitle{特征函数的定义}
	\begin{defi}
		设 $X$ 是一个随机变量，称
		\begin{eqnarray*}
			\varphi(t)&:=&E(e^{itX})=\pause E[\cos(tX)]+iE[\sin(tX)]\\
			&=&\pause \int_{-\infty}^{+\infty}\cos(tx)dF(x)+i\int_{-\infty}^{+\infty}\sin(tx)dF(x)\\
			&=&\pause \int_{-\infty}^{+\infty}e^{itx}dF(x), \quad -\infty <t<+\infty
		\end{eqnarray*}
		\pause 为 $X$ 的特征函数。特别的，
		\begin{itemize}[<+-|alert@+>]
			\item
			如果 $X$ 为离散型随机变量，其分布列为 $p_k=P (X=x_k)$, 则
			\begin{eqnarray*}
				\varphi(t)=\sum_{k=1}^{+\infty}e^{itx_k}p_k, \quad -\infty <t<+\infty;
			\end{eqnarray*}
			\item  如果 $X$ 为连续型随机变量，其分布密度为 $p (x)$, 则
			\begin{eqnarray*}
				\varphi(t)=\int_{-\infty}^{+\infty}e^{itx}p(x)dx, \quad -\infty <t<+\infty;
			\end{eqnarray*}
		\end{itemize}

	\end{defi}
\end{frame}


\begin{frame}%[allowframebreaks]
	\frametitle{特征函数的性质 I}
	\begin{enumerate}[<+-|alert@+>]
		\item  由 $\varphi (t)=E[\cos (tX)]+iE[\sin (tX)]$ 知，特征函数总是存在的；
		\item $\varphi (0)=1, |\varphi (t)|\le 1$: $\varphi (0)=E (1)=1$, 而
		\begin{eqnarray*}
			|\varphi(t)|&=&|E[\cos(tX)]+iE[\sin(tX)]|=\pause \big[(E[\cos(tX)])^2+(E[\sin(tX)])^2\big]^{1/2}\\
			&\le&\pause \big[E[\cos^2(tX)]+E[\sin^2(tX)]\big]^{1/2}=\pause 1;
		\end{eqnarray*}
		\item $\varphi(-t)=\overline{\varphi(t)}$: $\varphi(-t)=E[\cos(tX)]-iE[\sin(tX)]=\overline{\varphi(t)}$;
		\item  若 $Y=aX+b$, 其中 $a,b$ 为常数，则 $\varphi_Y (t)=e^{ibt}\varphi_X (at)$, 事实上，
		{\small\begin{eqnarray*}
				&&\hspace{-1.2cm}\varphi_Y(t)=\pause E[e^{it(aX+b)}]=E[\cos(t(aX+b))]+iE[\sin(t(aX+b))]\\
				&&\hspace{-1.2cm}=\pause E[\cos(taX)\cos(tb)-\sin(taX)\sin(tb)]+iE[\sin(taX)\cos(tb)+\cos(taX)\sin(tb)]\\
				&&\hspace{-1.2cm}=\pause \cos(tb)\big(E[\cos(taX)]+iE[\sin(taX)]\big)+i\sin(tb)\big(E[\cos(taX)]+iE[\sin(taX)]\big)\\
				&&\hspace{-1.2cm}=\pause  e^{itb}Ee^{itaX}=e^{itb}\varphi_X(at)\pause
		\end{eqnarray*}}
	\end{enumerate}
\end{frame}
\begin{frame}
	\frametitle{特征函数的性质 II}

	\begin{enumerate}[<+-|alert@+>]
		\item[5.] 若 $X,Y$ 独立，则 $\varphi_{X+Y}(t)=\varphi_X (t)\varphi_Y (t)$, 事实上，
		\begin{eqnarray*}
			&&\varphi_{X+Y}(t)=E(e^{it(X+Y)})=\pause E(e^{itX}e^{itY})\\
			&&=\pause E[(\cos(tX)+i\sin(tX))(\cos(tY)+i\sin(tY))]\\
			&&=\pause E[(\cos(tX)+i\sin(tX))\cos(tY)]+iE[(\cos(tX)+i\sin(tX))\sin(tY)]\\
			&&=\pause E[\cos(tY)]E[e^{itX}]+iE[\sin(tY)]E[e^{itX}]\\
			&&=\pause E(e^{itX})E(e^{itY})
		\end{eqnarray*}\pause
		\item[6.] 更一般的，若 $X_1,X_2,\cdots,X_n$ 相互独立，则对于 $Y=X_1+\cdots+X_n$ 的特征函数有
		\begin{eqnarray*}
			\varphi_Y(t)=\Pi_{k=1}^n\varphi_{X_k}(t)
		\end{eqnarray*}
	\end{enumerate}
\end{frame}
\begin{frame}
	\frametitle{特征函数的性质 III}

	\begin{enumerate}
		\item[7.] 如果 $E (X^k)$ 存在，则
		\begin{eqnarray*}
			\varphi^{(k)}(t)=i^kE[X^ke^{itX}], \quad     \varphi^{(k)}(0)=i^kE[X^k]
		\end{eqnarray*}
		事实上，
		\begin{eqnarray*}
			\dfrac{d^k}{dt^k}\varphi(t)=\pause  \dfrac{d^k}{dt^k}E[e^{itX}]=\pause E[\dfrac{d^k}{dt^k}e^{itX}]=\pause E[(iX)^ke^{itX}]=\pause i^kE[X^ke^{itX}]\pause
		\end{eqnarray*}
	\end{enumerate}
\end{frame}
\begin{frame}
	\frametitle{特征函数的性质 IV}

	\begin{enumerate}
		\item[8.] $\varphi (t)$ 在 $(-\infty,+\infty)$ 上一致连续:
		\begin{eqnarray*}
			&&|\varphi(t+h)-\varphi(t)|=\pause |E(e^{i(t+h)X})-E(e^{itX})|=\pause |E[e^{itX}(e^{ihX}-1)]|\\
			&&\le\pause E|e^{itX}(e^{ihX}-1)|=\pause E|e^{ihX}-1|=\pause \int_{-\infty}^{+\infty}|e^{ihx}-1|dF_X(x)\\
			&&=\pause \int_{|x|>A}|e^{ihx}-1|dF_X(x)+\int_{-A}^A|e^{ihx}-1|dF_X(x)\\
			&&\le\pause 2 \int_{|x|>A}dF_X(x)+\int_{-A}^A|hx|dF_X(x)\\
			&&\le\pause 2P(|X|>A)+hAP(|X|\le A)\le \pause 2P(|X|>A)+hA
		\end{eqnarray*}
		\pause 选取 $A$ 使得 $P (|X|>A)\le \epsilon/4$, $h\le \delta=\dfrac{\epsilon}{2A}$, 则只要 $|h|\le \delta$ 必有
		\pause \begin{eqnarray*}
			|\varphi(t+h)-\varphi(t)|\le 2P(|X|>A)+hA\le \epsilon
		\end{eqnarray*}
	\end{enumerate}
\end{frame}
\begin{frame}
	\frametitle{特征函数的性质 V}

	\begin{enumerate}
		\item[9.] $\varphi (t)$ 是非负定函数，即对任意的正整数 $n$ 及 $n$ 个实数 $t_1,\cdots, t_n$ 及 $n$ 个复数 $z_1,\cdots, z_n$, 均有
		\begin{eqnarray*}
			\sum_{k=1}^n\sum_{j=1}^n\varphi(t_k-t_j)z_k\bar{z}_j\ge 0
		\end{eqnarray*}
		\zheng 由特征函数的定义可得
		{\small\begin{eqnarray*}
				\sum_{k=1}^n\sum_{j=1}^n\varphi(t_k-t_j)z_k\bar{z}_j&=&\pause \sum_{k=1}^n\sum_{j=1}^nz_k\bar{z}_jE[e^{i(t_k-t_j)X}]=\sum_{k=1}^n\sum_{j=1}^nE[z_ke^{it_kX}\bar{z}_j\overline{e^{it_jX}}]\\
				&=&\pause E[\sum_{k=1}^nz_ke^{it_kX}\sum_{j=1}^n\overline{z_je^{it_jX}}]=\pause E\big[|\sum_{k=1}^nz_ke^{it_kX}|^2\big]\ge 0
		\end{eqnarray*}}

	\end{enumerate}
\end{frame}



\begin{frame}%[allowframebreaks]
	\frametitle{常用分布的特征函数  I}
	\begin{itemize}[<+-|alert@+>]
		\item 单点分布 $P (X=a)=1$: $\varphi (t)=e^{ita}$;
		\item 两点分布  $P (X=x)=p^x (1-p)^{1-x}, x=0,1$: $\varphi (t)=pe^{it}+(1-p)$;
		\item 泊松分布:$\varphi (t)=\sum_{k=0}^{+\infty} e^{ikt}\dfrac{\lambda^k}{k!} e^{-\lambda}=e^{\lambda e^{it}} e^{-\lambda}=e^{\lambda (e^{it}-1)}$;
		\item 均匀分布 $U (a,b)$:$\varphi (t)=\int_a^b\dfrac{e^{itx}}{b-a} dx=\dfrac{e^{ibt}-e^{iat}}{it (b-a)}$;
		\item 标准正态分布:
		\begin{eqnarray*}
			\hspace{-1cm}\varphi(t)&=&\pause \int_{-\infty}^{+\infty}e^{itx}\dfrac{1}{\sqrt{2\pi}}e^{-\frac{x^2}{2}}dx\\
			&=&\pause \dfrac{1}{\sqrt{2\pi}}\bigg(\int_{-\infty}^{+\infty}\cos(tx) e^{-\frac{x^2}{2}}dx+i\int_{-\infty}^{+\infty}\sin(tx) e^{-\frac{x^2}{2}}dx\bigg)\\
			\hspace{-1cm}&=&\pause \dfrac{1}{\sqrt{2\pi}}\int_{-\infty}^{+\infty}\cos(tx) e^{-\frac{x^2}{2}}dx
		\end{eqnarray*}\pause
	\end{itemize}
\end{frame}
\begin{frame}%[allowframebreaks]
	\frametitle{常用分布的特征函数  II}
	故 \begin{eqnarray*}
		\varphi'(t)&=&\pause -\dfrac{1}{\sqrt{2\pi}}\int_{-\infty}^{+\infty}x\sin(tx) e^{-\frac{x^2}{2}}dx=\pause \dfrac{1}{\sqrt{2\pi}}\int_{-\infty}^{+\infty}\sin(tx) de^{-\frac{x^2}{2}}\\
		&=&\pause -\dfrac{1}{\sqrt{2\pi}}\int_{-\infty}^{+\infty}t\cos(tx) e^{-\frac{x^2}{2}}dx=\pause -t\varphi(t)
	\end{eqnarray*}
	\pause 从而
	\begin{eqnarray*}
		\dfrac{d}{dt}[\varphi(t)e^{t^2/2}]=[\varphi'(t)+t\varphi(t)]e^{t^2/2}=0
	\end{eqnarray*}
	\pause 故 $\varphi (t) e^{t^2/2}=C$, 再利用 $C=\varphi (0)=1$ 得
	\begin{eqnarray*}
		\varphi(t)=e^{-t^2/2}.\pause
	\end{eqnarray*}
\end{frame}
\begin{frame}%[allowframebreaks]
	\frametitle{常用分布的特征函数  III}
	\begin{itemize}[<+-|alert@+>]
		\item 指数分布:
		\begin{eqnarray*}
			\varphi(t)=\int_0^\infty e^{itx}\lambda e^{-\lambda x}dx=\bigg(1-\dfrac{it}{\lambda}\bigg)^{-1}
		\end{eqnarray*}
		\item 二项分布: $Y\sim B (n,p)$. 注意到 $Y=X_1+\cdots+X_n$, 其中 $X_i$ 为独立同分布的随机变量，且 $X_i\sim B (1,p)$, 故
		\[\varphi_Y(t)=\Pi_{k=1}^n(pe^{it}+q)=(pe^{it}+q)^n, \quad q=1-p\]

		\item 一般正态分布: $Y\sim N (\mu,\sigma^2)$, 则由 $X:=\dfrac{Y-\mu}{\sigma}\sim N (0,1)$ 及 $Y=\sigma X+\mu$ 可得
		\begin{eqnarray*}
			\varphi_Y(t)=e^{i\mu t}\varphi_X(\sigma t)=e^{i\mu t-\frac{\sigma^2t^2}{2}}
		\end{eqnarray*}
	\end{itemize}
\end{frame}
\begin{frame}%[allowframebreaks]
	\frametitle{常用分布的特征函数  IV}
	\begin{itemize}[<+-|alert@+>]
		\item 伽玛分布:$Y\sim \Gamma (n,\lambda)$. 注意到 $Y=X_1+\cdots+X_n$, 其中 $X_i$ 独立同指数分布即 $X_i\sim \exp (\lambda)$, 故
		\begin{eqnarray*}
			\varphi_Y(t)=\Pi_{k=1}^n(1-\dfrac{it}{\lambda})^{-1}=(1-\dfrac{it}{\lambda})^{-n}
		\end{eqnarray*}
		更进一步，当 $Y\sim \Gamma (\alpha,\lambda)$, 我们有
		\begin{eqnarray*}
			\varphi_Y(t)=(1-\dfrac{it}{\lambda})^{-\alpha}
		\end{eqnarray*}
		\item $\chi^2 (n)$ 分布：因为 $\chi^2 (n)=\Gamma (n/2,1/2)$, 故其特征函数为
		\begin{eqnarray*}
			\varphi(t)=(1-2it)^{-n/2}
		\end{eqnarray*}
	\end{itemize}
\end{frame}
\begin{frame}
	\begin{exam}
		利用特征函数的方法求 $\Gamma (\alpha,\lambda)$ 分布的数学期望与方差.
	\end{exam}

	\pause
	\jieda $\Gamma$ 分布的特征函数及其导数分别为
	\begin{eqnarray*}
		\varphi(t)&=&\pause \big(1-\dfrac{it}{\lambda}\big)^{-\alpha}\\
		\varphi'(t)&=&\pause \dfrac{i\alpha}{\lambda}\big(1-\dfrac{it}{\lambda}\big)^{-\alpha-1},\quad \pause \varphi'(0)=\dfrac{i\alpha}{\lambda}\\
		\varphi''(t)&=&\pause \dfrac{i^2\alpha(\alpha+1)}{\lambda^2}\big(1-\dfrac{it}{\lambda}\big)^{-\alpha-2}, \quad \pause \varphi''(0)=-\dfrac{\alpha(\alpha+1)}{\lambda^2}
	\end{eqnarray*}\pause
	故由 $\varphi^{(k)}(0)=i^kE (X^k)$ 知，\pause
	\begin{eqnarray*}
		E(X)&=&\dfrac{\varphi'(0)}{i}=\dfrac{\alpha}{\lambda},\quad E(X^2)=\dfrac{\varphi''(0)}{i^2}=\dfrac{\alpha(\alpha+1)}{\lambda^2}\\
		D(X)&=&\pause E(X^2)-(EX)^2=\pause \dfrac{\alpha(\alpha+1)}{\lambda^2}-\dfrac{\alpha^2}{\lambda^2}=\pause \dfrac{\alpha}{\lambda^2}
	\end{eqnarray*}

\end{frame}

\begin{frame}
	\frametitle{特征函数与分布函数之间的关系：相互唯一决定}
	\begin{thm}[逆转公式] 设 $F (x)$ 与 $\varphi (t)$ 分别为随机变量 $X$ 的分布函数和特征函数，则对 $F (x)$ 的任意两个连续点 $x_1<x_2$ 有
		\begin{eqnarray*}
			F(x_2)-F(x_1)=\lim_{T\rightarrow\infty}\dfrac{1}{2\pi}\int_{-T}^T\dfrac{e^{-itx_1}-e^{-itx_2}}{it}\varphi(t)dt
		\end{eqnarray*}
	\end{thm}
	\pause
	\begin{thm}[唯一性定理] 随机变量的分布函数由其特征函数唯一决定.
	\end{thm}

	\pause \zheng 在逆转公式中令 $x_2=x, x_1=y\rightarrow -\infty$ 可得
	\begin{eqnarray*}
		F(x)=\lim_{y\rightarrow-\infty}\lim_{T\rightarrow\infty}\dfrac{1}{2\pi}\int_{-T}^T\dfrac{e^{-ity}-e^{-itx}}{it}\varphi(t)dt
	\end{eqnarray*}
	\pause 而由分布函数本身是右连续可知，分布函数任一点的取值可由连续点上的值唯一决定.

\end{frame}

\begin{frame}%[allowframebreaks]
	\frametitle{狄利克雷积分 I}
	\begin{lem}[狄利克雷积分] 令
		\[I(a,T)=\int_0^T \dfrac{\sin at}{t}dt,\quad{\rm sgn}\{a\}:=\left\{
		\begin{array}{rl}
			-1,& a<0,\\
			0,& a=0,\\
			1,& a>0,
		\end{array}\right.\] 则
		\[\lim_{T\rightarrow+\infty}I(a,T)=\lim_{T\rightarrow+\infty}\int_0^T \dfrac{\sin at}{t}dt=\dfrac{\pi}{2}
		{\rm sgn}\{a\}\]
	\end{lem}

	\zheng 注意到如果作积分换元 $y=at$ 可得
	\begin{eqnarray*}
		\lim_{T\rightarrow+\infty}I(a,T)={\rm sgn}\{a\}\lim_{T\rightarrow+\infty}I(1,T)
	\end{eqnarray*}
\end{frame}
\begin{frame}%[allowframebreaks]
	\frametitle{狄利克雷积分 II}
	只需证明
	\begin{eqnarray*}
		\lim_{T\rightarrow+\infty}I(1,T)=\lim_{T\rightarrow+\infty}\int_0^T\dfrac{\sin t}{t}dt=\dfrac{\pi}{2}.
	\end{eqnarray*}
	注意到 $\dfrac{1}{t}=\int_0^{+\infty} e^{-ut} du$, 故
	\begin{eqnarray*}
		I(1,T)&=&\pause \int_0^T\int_0^{+\infty}e^{-ut}\sin tdu dt=\pause \int_0^{+\infty} \bigg[\int_0^Te^{-ut}\sin tdt\bigg] du\\
		&=&\pause  \int_0^{+\infty}\bigg[\dfrac{1}{1+u^2}-\dfrac{1}{1+u^2}e^{-uT}(u\sin T+\cos T)\bigg]du \\
		&=&\pause \dfrac{\pi}{2}- \int_0^{+\infty}\dfrac{u\sin T+\cos T}{1+u^2}e^{-uT}du\rightarrow \pause \dfrac{\pi}{2}
	\end{eqnarray*}

\end{frame}
\begin{frame}
	\frametitle{逆转公式的证明}
	\vspace{-0.2cm}
	\zheng 考虑积分
	\begin{eqnarray*}
		J(T)&=&\pause \dfrac{1}{2\pi}\int_{-T}^T\dfrac{e^{-itx_1}-e^{-itx_2}}{it}\varphi(t)dt\\
		&=&\pause \dfrac{1}{2\pi}\int_{-T}^T\bigg[\int_{-\infty}^{+\infty}\dfrac{e^{-itx_1}-e^{-itx_2}}{it}e^{itx}dF(x)\bigg]dt\\
		&=&\pause \dfrac{1}{2\pi}\int_{-\infty}^{+\infty}\bigg[\int_{-T}^T\dfrac{e^{it(x-x_1)}-e^{it(x-x_2)}}{it}dt\bigg]dF(x)\\
		&=&\pause \dfrac{1}{\pi}\int_{-\infty}^{+\infty}\bigg[\int_0^T\dfrac{\sin t(x-x_1)}{t}dt-\int_0^T\dfrac{\sin t(x-x_2)}{t}dt\bigg]dF(x)\\
		&=&\pause \dfrac{1}{\pi}\int_{-\infty}^{+\infty}\big[I(x-x_1,T)-I(x-x_2,T)\big]dF(x)\\
		&\rightarrow&\pause \dfrac{1}{\pi}\dfrac{\pi}{2}\int_{-\infty}^{+\infty}\big[{\rm sgn}\{x-x_1\}-{\rm sgn}\{x-x_2\}\big]dF(x) \qquad T\rightarrow+\infty\\
		&=&\pause     \int_{x_1}^{x_2}dF(x)=F(x_2)-F(x_1)
	\end{eqnarray*}

\end{frame}
\begin{frame}
	\frametitle{逆转公式的一般情形}
	\begin{thm}[逆转公式一般情形] 设 $F (x)$ 与 $\varphi (t)$ 分别为随机变量 $X$ 的分布函数和特征函数，则对任意的 $x_1,x_2\in R$ 有
		{\small \begin{eqnarray*}
				\dfrac{F(x_2-0)+F(x_2)}{2}-\dfrac{F(x_1-0)+F(x_1)}{2}=\lim_{T\rightarrow\infty}\dfrac{1}{2\pi}\int_{-T}^T\dfrac{e^{-itx_1}-e^{-itx_2}}{it}\varphi(t)dt
		\end{eqnarray*}}
	\end{thm}
\end{frame}
\begin{frame}
	\frametitle{连续型分布情形下的逆转公式}
	\begin{thm}
		若 $X$ 为连续型随机变量，其密度函数为 $p (x)$, 其特征函数为
		\[\varphi(t)=\int_{-\infty}^{+\infty}e^{itx}p(x)dx,\]
		如果 $\int_{-\infty}^{+\infty}|\varphi (t)|dt<+\infty$, 则
		\begin{eqnarray*}
			p(x)=\dfrac{1}{2\pi}\int_{-\infty}^{+\infty}e^{-itx}\varphi(t)dt
		\end{eqnarray*}
	\end{thm}
	\pause \zheng 由分布密度的定义可知
	\begin{eqnarray*}
		p(x)&=&\pause \lim_{h\rightarrow 0}\dfrac{F(x+h)-F(x)}{h}=\pause \lim_{h\rightarrow 0}\dfrac{1}{2\pi}\int_{-\infty}^{+\infty}\dfrac{e^{-itx}-e^{-it(x+h)}}{it\cdot h}\varphi(t)dt\\
		&=&\pause\dfrac{1}{2\pi}\int_{-\infty}^{+\infty} \lim_{h\rightarrow 0}\dfrac{e^{-ith}-1}{-it\cdot h}e^{-itx}\varphi(t)dt=\pause \dfrac{1}{2\pi}\int_{-\infty}^{+\infty}e^{-itx}\varphi(t)dt
	\end{eqnarray*}
\end{frame}
\end{document}
